% GENERATED FILE. Edit canonical lessons, then run npm run generate:books.
% canonical-source-hash: fnv1a64:dac3884559feae1a
% canonical-lessons: LA-C301-susurrus, LA-C301-susurro, LA-C301-vulgo, LA-C301-crebresco, LA-C301-perhibeo

\chapter{Rūmor et susurrus}
\label{ch:la-rumor-et-susurrus}
\hlchaptermodality{301}

\begin{chapteropening}
\textbf{By the end of this chapter:} \emph{I can tell how a rumour spread: a whisper, to whisper, to spread abroad, to grow stronger and they say.}

Complete the last lesson of chapter 301: I can tell how a rumour spread: a whisper, to whisper, to spread abroad, to grow stronger and they say.
\end{chapteropening}

\section[susurrus, susurrī]{susurrus, susurrī — a whisper, a murmur}

\label{lesson:LA-C301-susurrus}



\begin{quote}
Before the new one: say the Latin for an embassy, then the Latin for to bring back word.
\end{quote}

\subsection*{You'll want to know: susurrus, susurrī}
\textbf{susurrus, susurrī} — \textquotedblleft{}a whisper, a murmur\textquotedblright{}.

\begin{quote}
Per tōtam urbem susurrus erat.
\end{quote}

\begin{quote}
\emph{There was a whisper all through the city.}
\end{quote}

Masculine. A sound word: the hum of bees, the rustle of wind, the murmur of gossip. English \emph{susurration}.

\begin{cousinweb}
\textbf{susurrus} is a doubled word for a low hum, which Lewis \& Short build on a root \emph{sur} for sound. English \emph{susurration} comes from its family.
\end{cousinweb}

\begin{grammarlens}[title={What you've built}]
One more word for telling, explaining and describing at B1.
\end{grammarlens}

\subsection*{Your turn}
\begin{itemize}
\raggedright
  \item \emph{Say it:} \emph{susurrus, susurrī}
  \item \emph{Say it:} \emph{susurrus, susurrī}, once more
  \item \emph{Say it:} a sentence of your own with \emph{susurrus, susurrī}
  \item \emph{Recall:} say the Latin for an embassy, then the Latin for to bring back word, then say \emph{susurrus, susurrī} again
\end{itemize}

\begin{tcolorbox}[breakable,colback=teal!4,colframe=teal!35!black,title={Before you move on}]
What does susurrus, susurrī mean? (\textquotedblleft{}A whisper, a murmur\textquotedblright{}.) Say it once more.
\end{tcolorbox}
\section[susurrō, susurrāre]{susurrō, susurrāre — to whisper}

\label{lesson:LA-C301-susurro}



\begin{quote}
Before the new one: say the Latin for to bring back word, then the Latin for a whisper.
\end{quote}

\subsection*{You'll want to know: susurrō, susurrāre}
\textbf{susurrō, susurrāre} — \textquotedblleft{}to whisper\textquotedblright{}.

\begin{quote}
Ancilla dominae aliquid in aurem susurrāvit.
\end{quote}

\begin{quote}
\emph{The maid whispered something into her mistress's ear.}
\end{quote}

\emph{susurrō, susurrāre, susurrāvī, susurrātum}, mostly poetic (Vergil, Ovid). \emph{In aurem}: into the ear.

\begin{cousinweb}
\textbf{susurrō} is built on \emph{susurrus}, \textquotedblleft{}a whisper, a hum\textquotedblright{} (EtymDB), itself a sound word.
\end{cousinweb}

\begin{grammarlens}[title={What you've built}]
One more word for telling, explaining and describing at B1.
\end{grammarlens}

\subsection*{Your turn}
\begin{itemize}
\raggedright
  \item \emph{Say it:} \emph{susurrō, susurrāre}
  \item \emph{Say it:} \emph{susurrō, susurrāre}, once more
  \item \emph{Say it:} a sentence of your own with \emph{susurrō, susurrāre}
  \item \emph{Recall:} say the Latin for to bring back word, then the Latin for a whisper, then say \emph{susurrō, susurrāre} again
\end{itemize}

\begin{tcolorbox}[breakable,colback=teal!4,colframe=teal!35!black,title={Before you move on}]
What does susurrō, susurrāre mean? (\textquotedblleft{}To whisper\textquotedblright{}.) Say it once more.
\end{tcolorbox}
\section[vulgō, vulgāre]{vulgō, vulgāre — to spread abroad, to make public}

\label{lesson:LA-C301-vulgo}



\begin{quote}
Before the new one: say the Latin for a whisper, then the Latin for to whisper.
\end{quote}

\subsection*{You'll want to know: vulgō, vulgāre}
\textbf{vulgō, vulgāre} — \textquotedblleft{}to spread abroad, to make public\textquotedblright{}.

\begin{quote}
Rūmor dē pāce celeriter vulgātus est.
\end{quote}

\begin{quote}
\emph{The rumour about peace was quickly spread abroad.}
\end{quote}

\emph{vulgō, vulgāre, vulgāvī, vulgātum}, also written \emph{volgō}. Its compound \emph{dīvulgō} gives English \emph{divulge}.

\begin{cousinweb}
\textbf{vulgō} is built on \emph{vulgus}, \textquotedblleft{}the common people, the crowd\textquotedblright{} (EtymDB): to put something among the crowd. English \emph{vulgar} comes from \emph{vulgus}.
\end{cousinweb}

\begin{grammarlens}[title={What you've built}]
One more word for telling, explaining and describing at B1.
\end{grammarlens}

\subsection*{Your turn}
\begin{itemize}
\raggedright
  \item \emph{Say it:} \emph{vulgō, vulgāre}
  \item \emph{Say it:} \emph{vulgō, vulgāre}, once more
  \item \emph{Say it:} a sentence of your own with \emph{vulgō, vulgāre}
  \item \emph{Recall:} say the Latin for a whisper, then the Latin for to whisper, then say \emph{vulgō, vulgāre} again
\end{itemize}

\begin{tcolorbox}[breakable,colback=teal!4,colframe=teal!35!black,title={Before you move on}]
What does vulgō, vulgāre mean? (\textquotedblleft{}To spread abroad, to make public\textquotedblright{}.) Say it once more.
\end{tcolorbox}
\section[crēbrēscō, crēbrēscere]{crēbrēscō, crēbrēscere — to grow stronger, to spread (of a rumour)}

\label{lesson:LA-C301-crebresco}



\begin{quote}
Before the new one: say the Latin for to whisper, then the Latin for to spread abroad.
\end{quote}

\subsection*{You'll want to know: crēbrēscō, crēbrēscere}
\textbf{crēbrēscō, crēbrēscere} — \textquotedblleft{}to grow stronger, to spread (of a rumour)\textquotedblright{}.

\begin{quote}
Fāma dē victōriā in diēs crēbrēscēbat.
\end{quote}

\begin{quote}
\emph{The report of the victory grew stronger day by day.}
\end{quote}

\emph{crēbrēscō, crēbrēscere, crēbruī}, from \emph{crēber} (taught). Tacitus uses it of a rumour that keeps growing.

\begin{cousinweb}
\textbf{crēbrēscō} is built on \emph{crēber}, \textquotedblleft{}frequent, close-packed\textquotedblright{}, with the ending \emph{-ēscō} for something that is becoming so (Lewis \& Short): to become frequent.
\end{cousinweb}

\begin{grammarlens}[title={What you've built}]
One more word for telling, explaining and describing at B1.
\end{grammarlens}

\subsection*{Your turn}
\begin{itemize}
\raggedright
  \item \emph{Say it:} \emph{crēbrēscō, crēbrēscere}
  \item \emph{Say it:} \emph{crēbrēscō, crēbrēscere}, once more
  \item \emph{Say it:} a sentence of your own with \emph{crēbrēscō, crēbrēscere}
  \item \emph{Recall:} say the Latin for to whisper, then the Latin for to spread abroad, then say \emph{crēbrēscō, crēbrēscere} again
\end{itemize}

\begin{tcolorbox}[breakable,colback=teal!4,colframe=teal!35!black,title={Before you move on}]
What does crēbrēscō, crēbrēscere mean? (\textquotedblleft{}To grow stronger, to spread (of a rumour)\textquotedblright{}.) Say it once more.
\end{tcolorbox}
\section[perhibeō, perhibēre]{perhibeō, perhibēre — to assert, to say (that)}

\label{lesson:LA-C301-perhibeo}



\begin{quote}
Before the new one: say the Latin for to spread abroad, then the Latin for to grow stronger.
\end{quote}

\subsection*{You'll want to know: perhibeō, perhibēre}
\textbf{perhibeō, perhibēre} — \textquotedblleft{}to assert, to say (that)\textquotedblright{}.

\begin{quote}
Hunc fluvium altissimum esse perhibent.
\end{quote}

\begin{quote}
\emph{They say this river is very deep.}
\end{quote}

\emph{perhibeō, perhibēre, perhibuī, perhibitum}: \emph{per} + \emph{habeō}. \emph{Perhibent} + accusative and infinitive: 'they say that…', the way rumour talks.

\begin{cousinweb}
\textbf{perhibeō} is built on \emph{habeō}, \textquotedblleft{}to hold\textquotedblright{}, with \emph{per-} (EtymDB): to hold out, to put forward, and so to say.
\end{cousinweb}

\begin{grammarlens}[title={What you've built}]
One more word for telling, explaining and describing at B1.
\end{grammarlens}

\subsection*{Your turn}
\begin{itemize}
\raggedright
  \item \emph{Say it:} \emph{perhibeō, perhibēre}
  \item \emph{Say it:} \emph{perhibeō, perhibēre}, once more
  \item \emph{Say it:} a sentence of your own with \emph{perhibeō, perhibēre}
  \item \emph{Recall:} say the Latin for to spread abroad, then the Latin for to grow stronger, then say \emph{perhibeō, perhibēre} again
\end{itemize}

\begin{tcolorbox}[breakable,colback=teal!4,colframe=teal!35!black,title={Before you move on}]
What does perhibeō, perhibēre mean? (\textquotedblleft{}To assert, to say (that)\textquotedblright{}.) Say it once more.
\end{tcolorbox}
